\begin{partbacktext}
\part{Applications}
\noindent
The three chapters of this part apply the singularity theory developed earlier to three rather different settings: big toric geometry, non-Archimedean pluripotential theory and partial Bergman kernels. They are largely independent of one another, and each uses a different part of the machinery built in the preceding chapters.

\cref{chap:toric_big} extends the toric dictionary of \cref{chap:toric_ample} from ample line bundles to big toric line bundles. In the big case, identifying the Newton body is no longer as simple as in the ample setting, and the behavior of potentials along toric strata becomes more delicate. The chapter establishes the toric-convex dictionary using the theory of partial Okounkov bodies developed in \cref{chap:Okounkov}.

\cref{chap:NApotential} develops a non-Archimedean pluripotential theory on compact K\"ahler manifolds. The basic objects are inverse systems of $\mathcal I$-model test curves, and the operations on test curves from \cref{sec:operationtc} provide the corresponding algebra of non-Archimedean potentials. The chapter constructs Duistermaat--Heckman measures and proves a comparison theorem identifying the resulting space with the Boucksom--Jonsson space of non-Archimedean psh metrics in the projective setting.

\cref{chap:Bergman} proves the convergence of partial Bergman kernels. Given a singular $\theta$-psh potential, a closed non-pluripolar weighted subset $(K,v)$ and a Bernstein--Markov measure, the normalized partial Bergman measure converges weakly to the corresponding partial equilibrium measure. The proof combines the envelope theory of \cref{chap:envelopes} and the asymptotic Riemann--Roch theorem \cref{thm:DXmain1}. It extends the theorem of Berman--Boucksom--Witt Nystr\"om \cite{BBWN11} to the singular relative setting considered here.

The prerequisites for the three chapters are different. The big toric chapter uses basic toric geometry, for which \cite{CLS11} and \cite{Fultor} are standard references. The non-Archimedean chapter is best read with the Boucksom--Jonsson theory in mind, especially \cite{BJ22}. The Bergman kernel chapter is closest in spirit to \cite{BBWN11}. These references are not logically required for following the arguments, but they provide useful context.
\end{partbacktext}
